\section{点分治}\label{sec:point-divide}

\begin{algorithm}[H]
    \caption{点分治}
    \begin{algorithmic}[1]
        \Procedure{Solve}{$n$}
            \State \Call{Decompose}{1, n}
        \EndProcedure
        \Procedure{Decompose}{$u, \pvar{whole}$}
            \State $\pvar{rt} \gets 0$
            \State \Call{FindRoot}{u, \pvar{whole}, \pvar{rt}, 0}
            \State $\pvar{vis}[\pvar{rt}] \gets \mathrm{true}$
            \State \Call{Calc}{\pvar{rt}}
            \For{$v \in G[\pvar{rt}]$}
                \If{$\pvar{vis}[v]$}
                    \State \textbf{continue}
                \EndIf
                \State $\pvar{ss} \gets \begin{cases} \pvar{siz}[v], & \text{if }\,\pvar{siz}[v] < \pvar{siz}[\pvar{rt}], \\ \pvar{whole} - \pvar{siz}[\pvar{rt}], & \text{otherwise}. \end{cases}$
                \State \Call{Decompose}{v, \pvar{ss}}
            \EndFor
        \EndProcedure
        \Procedure{FindRoot}{$u, \pvar{whole}, \mathrm{out}\ \pvar{rt}, \pvar{fa}$}
            \State $\pvar{siz}[u] \gets 1$
            \State $\pvar{mxs}[u] \gets 0$
            \For{$v \in G[u]$}
                \If{$\pvar{vis}[v]$ \textbf{or} $v = \pvar{fa}$}
                    \State \textbf{continue}
                \EndIf
                \State \Call{FindRoot}{v, \pvar{whole}, \pvar{rt}, u}
                \State $\pvar{siz}[u] \gets \pvar{siz}[u] + \pvar{siz}[v]$
                \State $\pvar{mxs}[u] \gets \max(\pvar{mxs}[u], \pvar{siz}[v])$
            \EndFor
            \State $\pvar{mxs}[u] \gets \max(\pvar{mxs}[u], \pvar{whole} - \pvar{siz}[u])$
            \If{$2 \cdot \pvar{mxs}[u] \le \pvar{whole}$}
                \State $\pvar{rt} \gets u$
            \EndIf
        \EndProcedure
        \Procedure{Calc}{$\pvar{rt}$}
            \State \Comment{problem-specific; do not modify \pvar{vis}, \pvar{siz}, \pvar{mxs}}
        \EndProcedure
    \end{algorithmic}
\end{algorithm}

\subsection{待实现}

\code{Calc(rt)} 统计经过 \code{rt} 的路径。按子树分批处理：先查询当前结构，再插入该子树，避免同子树点对。

\subsection{数据契约}

\code{siz} 以本次 \code{FindRoot} 的入口 \code{u} 为根，故 \code{ss} 要按 \code{siz[v] < siz[rt]} 分儿子 / 父亲，父亲一侧取 \code{whole - siz[rt]}。

\code{Decompose} 在 \code{Calc} 之后才用 \code{siz} 算 \code{ss}，因此 \code{Calc} 不得改写 \code{vis}、\code{siz}、\code{mxs}。

\code{FindRoot} 递归深度为连通块高度，链上可达 $O(n)$；\code{Decompose} 的递归深度才是 $O(\log n)$。
